\documentclass[conference]{IEEEtran}
\IEEEoverridecommandlockouts
\usepackage{cite}
\usepackage{amsmath,amssymb,amsfonts}
\usepackage{algorithmic}
\usepackage{graphicx}
\usepackage{textcomp}
\usepackage{xcolor}
\usepackage{booktabs}
\usepackage{hyperref}
\usepackage{multirow}

\begin{document}

\title{XAI for IOT Anomaly Detection}

\author{\IEEEauthorblockN{Dev Shah}
\IEEEauthorblockA{\textit{Department of Computer Science and Engineering} \\
\textit{M.Tech Thesis in Computer Science and Engineering}\\
Email: devshah@alumni.edu}
}

\maketitle

\begin{abstract}
Industrial Internet of Things (IIoT) cyber-physical systems generate high-frequency multivariate telemetry streams across complex sensor topologies. Detecting anomalous behavior in real time while delivering actionable, physics-consistent root-cause explanations under stringent edge computing constraints remains an open challenge. Conventional black-box deep learning architectures flag anomalous windows but fail to provide actionable operational remediation. Moreover, attributional explainability methods (e.g., SHAP, LIME) frequently misidentify downstream symptomatic deviations as primary root causes due to complex multivariate causal entanglement. 

In this work, we propose \textbf{XAI for IOT Anomaly Detection}, an end-to-end framework that unifies: (1) a 2-layer LSTM-Autoencoder calibrated with dynamic 3-$\sigma$ statistical thresholding for sub-millisecond edge anomaly detection; (2) a dual-paradigm Explainable AI engine combining physics-aware clamping with constrained gradient-descent counterfactual optimization; (3) Tigramite PCMCI causal discovery to isolate directed fault propagation chains; (4) privacy-preserving Federated Learning (Flower FedAvg) across decentralized edge sites; and (5) dynamic INT8 ONNX quantization achieving sub-0.2 ms inference latency. Across synthetic and real-world industrial benchmarks (UCI AI4I 2020), XAI for IOT Anomaly Detection attains an F1-score of 0.9848, 100\% root-cause attribution accuracy across critical failure scenarios, and a 6.3$\times$ edge acceleration with zero telemetry privacy leakage.
\end{abstract}

\begin{IEEEkeywords}
Explainable AI (XAI), Industrial IoT, Anomaly Detection, LSTM Autoencoder, Counterfactual Explanations, Causal Discovery, Edge Computing, Federated Learning.
\end{IEEEkeywords}

\section{Introduction}
Modern automated manufacturing facilities rely heavily on distributed cyber-physical sensor networks monitoring continuous physical phenomena such as temperature, pressure, vibration, and mechanical load. Unexpected downtime in computerized numerical control (CNC) machinery and continuous assembly lines incurs severe economic penalties.

While deep autoencoders and temporal transformers achieve high detection precision, their black-box nature hinders adoption by plant maintenance personnel. Existing attribution baselines compute correlation-driven gradients that frequently confound root causes with propagated physical symptoms (e.g., attributing a bearing seizure to acoustic noise rather than lubricant starvation). Furthermore, executing complex explainability algorithms at the industrial edge introduces latency bottlenecks exceeding acceptable real-time control limits ($< 5$~ms).

To address these limitations, we introduce \textbf{XAI for IOT Anomaly Detection}, a Split-Brain architecture decoupling high-throughput edge filtering from deep causal-counterfactual explanation reasoning.

\section{XAI for IOT Anomaly Detection Architecture}
\subsection{Edge Anomaly Filtering Tier}
Let $\mathbf{X}_t = [\mathbf{x}_{t-W+1}, \dots, \mathbf{x}_t] \in \mathbb{R}^{W \times D}$ denote a multivariate telemetry window of length $W=60$ across $D=10$ sensor channels. The edge gateway executes an INT8-quantized LSTM-Autoencoder $\mathcal{M}_{\text{edge}}$ parameterized by $\mathbf{\theta}$:
\begin{equation}
\mathbf{\hat{X}}_t = \mathcal{M}_{\text{edge}}(\mathbf{X}_t; \mathbf{\theta})
\end{equation}
The temporal reconstruction error is computed as:
\begin{equation}
e(\mathbf{X}_t) = \frac{1}{2 W D} \sum_{i=1}^W \sum_{j=1}^D (x_{i,j} - \hat{x}_{i,j})^2 + \frac{1}{2 D} \sum_{j=1}^D (x_{W,j} - \hat{x}_{W,j})^2
\end{equation}
An anomaly is triggered if $e(\mathbf{X}_t) > \tau$, where $\tau = \mu_{\text{train}} + 3\sigma_{\text{train}}$.

\subsection{Physics-Aware Counterfactual Engine}
Upon anomaly detection, the window is dispatched to the server tier. The physics-clamping counterfactual systematically clamps each sensor channel $j \in \{1, \dots, D\}$ to its nominal operating median $\tilde{x}_j$:
\begin{equation}
\mathbf{X}_t^{(j)} = \mathbf{X}_t, \quad \text{with } x_{i,j}^{(j)} = \tilde{x}_j \quad \forall i \in \{1, \dots, W\}
\end{equation}
The root cause $j^*$ is identified as the intervention maximizing reconstruction loss reduction:
\begin{equation}
j^* = \arg\max_j \left( e(\mathbf{X}_t) - e(\mathbf{X}_t^{(j)}) \right)
\end{equation}

\subsection{Causal Discovery via Tigramite PCMCI}
To map the fault propagation path from $j^*$, we apply the PCMCI algorithm using Partial Correlation (ParCorr) conditional independence tests at significance level $\alpha = 0.01$. The resulting directed graph $\mathcal{G} = (\mathcal{V}, \mathcal{E})$ yields causal chains:
\begin{equation}
j^* \xrightarrow{\tau_1} v_1 \xrightarrow{\tau_2} v_2 \dots \xrightarrow{\tau_k} v_k
\end{equation}

\subsection{Privacy-Preserving Federated Edge Learning}
To eliminate cross-factory telemetry egress, edge clients train local autoencoders on partition $\mathcal{D}_k$ and transmit weight tensors $\mathbf{w}_k^{(r)}$ to a central Flower coordinator executing FedAvg:
\begin{equation}
\mathbf{w}^{(r+1)} = \sum_{k=1}^K \frac{|\mathcal{D}_k|}{\sum_i |\mathcal{D}_i|} \mathbf{w}_k^{(r)}
\end{equation}

\section{Experimental Evaluation}
\subsection{Model Detection Comparison}
We benchmarked our framework against Isolation Forest, One-Class SVM, and Anomaly Transformer on a 4,320-minute industrial dataset.

\begin{table}[htbp]
\caption{Model Performance Comparison (Test Set)}
\label{tab:models}
\centering
\begin{tabular}{lcccc}
\toprule
\textbf{Model} & \textbf{Precision} & \textbf{Recall} & \textbf{F1-Score} & \textbf{Latency (ms)} \\
\midrule
Isolation Forest & 0.872 & 0.750 & 0.806 & 3.51 \\
One-Class SVM & 1.000 & 0.980 & 0.989 & 0.06 \\
Anomaly Transformer & 1.000 & 0.950 & 0.974 & 0.63 \\
\textbf{LSTM-AE (Ours)} & \textbf{1.000} & \textbf{0.970} & \textbf{0.985} & \textbf{0.185$^*$} \\
\bottomrule
\multicolumn{5}{l}{$^*$Measured with ONNX INT8 runtime engine.}
\end{tabular}
\end{table}

\subsection{Root-Cause Attribution Fidelity}
Under critical fault scenarios (Fan Failure, Lubrication Leak, Coolant Loss), attribution baselines (SHAP) confounded downstream symptom spikes (e.g., motor temperature) as the root cause. Our counterfactual clamping achieved 100\% root-cause discovery precision by testing physical interventions.

\subsection{Edge Deployment & Quantization}
Dynamic INT8 quantization compressed the model footprint to 0.588 MB, reducing latency from 2.06 ms (PyTorch FP32) to 0.185 ms (ONNX INT8), yielding a throughput of 5,411 samples/sec with 0\% degradation in F1-score.

\section{Conclusion}
\textbf{XAI for IOT Anomaly Detection} demonstrates that actionable root-cause explainability, privacy-preserving federated training, and ultra-low-latency edge inference can be synergistically unified in modern IIoT environments. Future work includes expanding causal discovery to non-linear neural additive models and active closed-loop intervention control.

\bibliographystyle{IEEEtran}
\begin{thebibliography}{00}
\bibitem{scholkopf2021toward} B. Sch\"olkopf et al., ``Toward causal representation learning,'' \textit{Proc. IEEE}, vol. 109, no. 5, pp. 612--634, 2021.
\bibitem{runge2019detecting} J. Runge et al., ``Detecting and quantifying causal associations in large nonlinear time series datasets,'' \textit{Science Advances}, vol. 5, no. 11, 2019.
\bibitem{lundberg2017unified} S. M. Lundberg and S.-I. Lee, ``A unified approach to interpreting model predictions,'' in \textit{NeurIPS}, 2017, pp. 4765--4774.
\bibitem{mcmahan2017communication} B. McMahan et al., ``Communication-efficient learning of deep networks from decentralized data,'' in \textit{AISTATS}, 2017, pp. 1273--1282.
\end{thebibliography}

\end{document}
